% Generated draft: requires mathematical and source review.
\documentclass[11pt]{article}
\usepackage[margin=25mm]{geometry}
\usepackage{fontspec}
\setmainfont{Latin Modern Roman}
\usepackage{amsmath,amssymb,mathtools,unicode-math}
\setmathfont{Latin Modern Math}
\usepackage{enumitem,needspace,xurl,hyperref}
\setlength{\parindent}{0pt}
\setlength{\parskip}{5pt}
\setlength{\emergencystretch}{3em}
\newcommand{\N}{\mathbb{N}}
\newcommand{\Z}{\mathbb{Z}}
\newcommand{\Q}{\mathbb{Q}}
\newcommand{\R}{\mathbb{R}}
\newcommand{\C}{\mathbb{C}}
\newcommand{\F}{\mathbb{F}}
\newcommand{\A}{\mathbb{A}}
\newcommand{\PP}{\mathbb{P}}
\newcommand{\E}{\mathbb{E}}
\newcommand{\eps}{\varepsilon}
\newcommand{\dd}{\,\mathrm{d}}
\begin{document}
\section*{UnsolvedMath ID: 30006237}
\textbf{Source problem code:} catalogue entry
\subsection{Definitions and mathematical statement}
Let \(p\) be a prime. The terms \emph{locality at \(p\)}, \emph{simple locality}, \emph{locally finite locality}, \emph{compact locality}, and the invariant \(\dim_p\) retain their intended meanings in the supplied report; their definitions are not included in the supplied record. In particular, no interpretation of \(\dim_p\) as a familiar rank or dimension is assumed.

A \emph{framing} of \(\mathcal L\), as described in the supplied excerpt, is a sequence
\[
\gamma=\bigl(\mathcal L_n\xrightarrow{\alpha_n}\mathcal L_{n+1}\bigr)_{n\geq 1}
\]
of injective homomorphisms of finite localities at \(p\), together with an identification
\[
\varinjlim_{n\geq 1}\mathcal L_n\cong\mathcal L
\]
in the category of partial groups. Thus there are compatible partial-group homomorphisms \(\iota_n:\mathcal L_n\to\mathcal L\), satisfying \(\iota_{n+1}\alpha_n=\iota_n\), with the following universal property: for every partial group \(\mathcal M\) and compatible family of partial-group homomorphisms \(f_n:\mathcal L_n\to\mathcal M\), there is a unique partial-group homomorphism \(f:\mathcal L\to\mathcal M\) such that \(f\iota_n=f_n\) for every \(n\geq1\). Indexing from \(1\) is only a convention.

\textbf{Statement to be reviewed (the supplied excerpt's Theorem A\(^{\prime}\)).} For every prime \(p\) and every simple, locally finite locality \(\mathcal L\) at \(p\), there exists a framing \(\gamma\) of \(\mathcal L\) such that every \(\mathcal L_n\) is simple and
\[
\dim_p(\mathcal L_n)=\dim_p(\mathcal L_m)
\qquad\text{for all }m,n\geq1.
\]

\textbf{Source qualifications.} The excerpt separately labels as a theorem of Gonzales the assertion that every compact locality at \(p\) admits a framing by finite localities at \(p\). This contextual assertion supplies neither simplicity nor constant dimension. Its attribution does not establish authorship of Theorem A\(^{\prime}\). The record labels the target open, whereas the excerpt calls it a theorem. Moreover, the supplied literature assessment says that the conjecture is missing after the words ``we conjecture'', although the supplied original statement continues with the two theorem-labelled assertions above. These discrepancies remain unresolved. This draft transcribes the explicit target without asserting its truth or its current open status; a self-contained formulation requires verification of the omitted terminology and any hypotheses attached to it.
\subsection{Short English statement}
The supplied target asks whether every simple, locally finite locality at a prime \(p\) is a direct limit, in the category of partial groups, of finite simple localities at the same prime whose \(p\)-dimensions are all equal. The specialized definitions and the conflict between the supplied open label, theorem label, and literature assessment require human review.
\subsection{Sources}
\begin{itemize}
\item Supplied local UnsolvedMath registry record, permanent ID \(30006237\), titled \emph{Constant-Dimension Framings of Simple Locally Finite Localities}; exact source problem code \texttt{catalogue entry}. Collection: Oberwolfach Reports Open Problems; proposed year: 2025. No proposer is named.
\item \emph{Finite Groups, Fusion Systems and Applications} (2025), Oberwolfach Reports, pp.\ 689--752, DOI \texttt{10.4171/owr/2025/16}. Supplied source and literature-source link, labelled ``original OWR report'': \url{https://doi.org/10.4171/owr/2025/16}. The supplied background describes Oberwolfach Reports as documenting mathematical workshops and their open-problem discussions. The report has not been independently consulted for this draft.
\item The supplied excerpt attributes the compact-locality framing theorem to Gonzales and labels the requested assertion Theorem A\(^{\prime}\); it gives no separate attribution for Theorem A\(^{\prime}\).
\item Literature assessment dated \texttt{2026-08-22}: supplied status ``open''. The assessment states that the record ends immediately after ``we conjecture'' and therefore omits the conjecture. This is supplied dated curation, not independently verified; it conflicts with the continuation present in the supplied original statement. The accompanying curation explicitly calls for independent verification of time-sensitive status and cited claims.
\item The registry supplies the statement-status label \texttt{corrected\_verified}, with the qualification that only formula delimiters, notation, grammar, and typography were repaired and that mathematical content and scope were preserved. This supplied label is not independent mathematical verification.
\end{itemize}
\end{document}
